\chapter*{KATA PENGANTAR}
\addcontentsline{toc}{chapter}{KATA PENGANTAR}

Puji syukur saya panjatkan kepada Tuhan Yang Maha Esa sehingga buku manual \textit{AI Music Detector} ini dapat diselesaikan. Buku ini disusun sebagai tugas mata kuliah Pemrograman~1 di Fakultas Teknik Universitas Negeri Semarang.

\textit{AI Music Detector} adalah aplikasi yang memperkirakan apakah sebuah rekaman musik dibuat manusia atau dihasilkan kecerdasan buatan (AI). Buku ini menjelaskan cara memasang dan memakai aplikasinya, cara kerja model di baliknya, metode yang dipakai pada setiap tahap, cara melatih ulang model, serta hasil pengujiannya. Seluruh kode sumber terbuka untuk umum di \url{\RepoURL}.

Saya mengucapkan terima kasih kepada \DosenPengampu{} selaku dosen pengampu mata kuliah Pemrograman~1 atas bimbingan dan kesempatan untuk mengerjakan proyek ini. Terima kasih juga kepada penyusun dataset SONICS, FMA, dan MUSEBENCH yang membuka data mereka untuk penelitian.

Saya menyadari aplikasi dan buku ini masih memiliki keterbatasan, terutama dalam mengenali generator musik AI yang belum pernah dilihat model. Kritik dan saran sangat saya harapkan agar proyek ini dapat terus diperbaiki. Semoga buku ini bermanfaat bagi pembaca.

\vspace{0.7cm}
\begin{flushright}
Semarang, 26 November 2026\\[1.6cm]
\NamaPenulis\\
NIM \NIMPenulis
\end{flushright}
